% ==============================================================================
% T0/FFGFT Theory: Document 241
% Title: The Two Layers – Ratios and Conversion.
%        A pedagogical introduction to the methodological architecture of FFGFT.
% Author: Johann Pascher
% Date: May 2026
% References: Documents 014 (nat./SI), 015 (Natural-units systematics),
%             044 (fine-structure convention), 052 (mass elimination),
%             133 (fractal correction), 134 (c as convention),
%             011 (fine-structure derivation), 205 (narrative),
%             159 (concert-hall / mode structure), 161 (Ising),
%             167 (LiNbO3), 173 (quantum dot), 185 (embedding cost),
%             230 (Hilbert-space bridge)
% ==============================================================================

\section*{Preface}

This document is for readers who want to enter the methodological
architecture of FFGFT without first working through its full
mathematical apparatus. It does not repeat any derivation; it explains
a reading order.

The content is not new. The central claim --- that FFGFT operates on
two cleanly separated layers --- already appears at various points in
the document series (above all in Docs.~014, 015, 044, and 134). But
there it stands as a technical clarification alongside other technical
remarks, not as an architectural principle with its own voice. That is
what this document tries to provide.

A first-time reader otherwise risks getting stuck at a point that
resolves itself the moment one recognizes the two layers. Whoever does
not recognize them sees in every fractal correction factor a free
parameter and in every natural constant an independent quantity ---
missing the actual point: that FFGFT operates with exactly one free
parameter $\xi_{0} = 4/30000$, and that every other quantity is either
a pure ratio or arises through a geometrically determined scale-bridge.

\section{The starting point: the constants problem}

Ordinary school physics presents a list of ``fundamental constants of
nature'':
\begin{itemize}\setlength\itemsep{0pt}
  \item the speed of light $c$,
  \item the reduced Planck constant $\hbar$,
  \item the gravitational constant $G$,
  \item the Boltzmann constant $k_B$,
  \item the electric constant $\varepsilon_0$,
  \item the fine-structure constant $\alpha \approx 1/137$,
  \item the elementary charge $e$,
\end{itemize}
and so on. Each has its own unit, its own value, each appears
independent. The student sooner or later asks: \emph{Why precisely
these values? Why not others?}

The usual answer, slightly abbreviated, runs: ``We don't know. That's
just how it is.''

This answer is intellectually unsatisfying. It suggests that nature is
equipped with a dozen independently adjustable knobs, and that humans
are to measure them off without ever being able to understand them.
Sommerfeld, Eddington, Dirac, and Feynman never accepted this; each
in his own way tried to see behind the wall of ``fundamental constants''.

FFGFT makes the same attempt, but on a different path. The path begins
with a seemingly harmless observation: many of these ``constants'' are
not really constants at all. They are \emph{exchange rates between
choices of scale}, not independent natural quantities. This becomes
visible only when one reduces consistently, in two steps.

\section{First step: $\hbar = c = G = k_B = 1$}

The first step is standard and has been common in theoretical physics
for decades: one works in \emph{natural units}, in which the constants
$\hbar$, $c$, $G$, and $k_B$ are simply set equal to one (Docs.~014,
015).

What at first appears to be notational convenience is in fact a
structural statement. In SI units, length (meter), time (second), and
mass (kilogram) have independent dimensions. This independence is,
however, a \emph{convention}, not a property of nature. Setting $c = 1$
makes length and time the same quantity. Setting $\hbar = 1$ makes
energy and inverse time the same quantity. Setting also $G = 1$ and
$k_B = 1$, \emph{all} mechanical, thermodynamic, and gravitational
quantities collapse onto a single dimension: energy.

\begin{center}
\begin{tabularx}{\linewidth}{Xl}
\toprule
Quantity & Dimension \\
\midrule
Energy, mass, temperature, frequency, momentum & $[E]$ \\
Length, time, wavelength & $[E]^{-1}$ \\
Force & $[E]^{2}$ \\
Velocity & dimensionless \\
Electric charge & dimensionless \\
\bottomrule
\end{tabularx}
\end{center}

This table (more fully in Doc.~015) says something simple but
important: \emph{the apparent variety of physical quantities is to a
large extent an artifact of the choice of scale}. There is only one
fundamental dimension, and all others are powers of it. The ``many
constants'' of school physics are no longer independent --- they are
conversion factors between different energy powers.

This step alone is not specific to FFGFT. It belongs to the standard
toolkit of high-energy physics. But it is the precondition for the
second, less familiar step.

\section{Second step: $\alpha = 1$}

Even in natural units with $\hbar = c = G = k_B = 1$, one quantity
seems to remain fundamental, unexplained, ``magical'': the
fine-structure constant
\[
  \alpha \;=\; \frac{e^{2}}{4\pi\varepsilon_{0}\hbar c} \;\approx\; \frac{1}{137{,}036}.
\]

Feynman called it ``the greatest mystery of physics'': a dimensionless
number that appears out of nowhere and yet governs the whole of
chemistry and atomic physics.

Here FFGFT takes a decisive second step: \emph{$\alpha$ may also be set
equal to one, in a consistently extended system of units}. This is not
a trick, not convention in the sense of mere convenience, but a
structural insight known in physics under the name
\textbf{Heaviside-Lorentz convention} and discussed at length in
Doc.~044:

\begin{itemize}\setlength\itemsep{0pt}
  \item In SI units: $\alpha = e^{2}/(4\pi\varepsilon_{0}\hbar c) \approx 1/137$.
  \item In Gaussian units: $\alpha = e^{2}/(\hbar c) \approx 1/137$.
  \item In Heaviside-Lorentz units (with additionally $4\pi\varepsilon_{0} = 1$
        and a suitable definition of charge), $\alpha = e^{2}/(4\pi)$,
        and in a further-extended system even $\alpha = 1$.
\end{itemize}

All three systems describe the same physics. The numerical value
$1/137{,}036$ is therefore \emph{not} a property of the world, but a
property of how we parametrize the world. What changes is only
\emph{where} the information sits --- in $\alpha$, in the charge $e$,
or in the electric constant $\varepsilon_{0}$.

Only after one has seen through this does the real point become
visible: \emph{exactly one dimensionless quantity remains that cannot
be removed by any choice of scale}. In FFGFT this quantity is the
geometric parameter
\[
  \xi_{0} \;=\; \frac{4}{30000} \;\approx\; 1{.}33 \times 10^{-4}.
\]

Everything else --- $c$, $\hbar$, $G$, $k_B$, $\alpha$, and even
particle masses, insofar as they are taken as pure ratios of one
another --- follows from $\xi_{0}$ alone, or is a scale convention.

\subsection*{Important clarification: ``set to 1'' means ``factored out''}

The values $c, \hbar, G, k_B, \alpha$ that have been ``set to one'' do
not disappear from physics. They are \emph{temporarily factored-out
scale factors}, not quantities rationalized away.

As long as one reckons with \emph{ratios} --- mass ratios, frequency
ratios, dimensionless couplings ---, these factors cancel out, and no
numerical substitution is needed. The answer is sharp without $c$ or
$\hbar$ ever appearing.

But the moment one wants to generate a \emph{projected measurement
value} in SI units --- the electron mass in kilograms, a frequency in
hertz, an energy in joules ---, the values that were set to one must
be \emph{re-inserted}, together with the fractal correction. They are
therefore conditional bridge factors: absent in the pure geometry,
present in the projection. This is not a contradiction but the clean
form of the architecture: whoever stays inside the geometry does not
need them; whoever wants to see a measurement value brings them back.

\section{The insight: two layers}

From these two steps follows the methodological architecture that
gives this document its name. FFGFT operates on \emph{two cleanly
separated layers}:

\begin{description}
  \item[Layer 1 --- the ratios.] Everything that can be expressed as a
    pure ratio between two quantities of the same kind --- mass ratios
    such as $m_{\mu}/m_{e}$, frequency ratios, winding numbers on the
    torus, dimensionless couplings, the parameter $\xi_{0}$ itself ---
    is \emph{unit-invariant}. These ratios are what is ``really there''
    independent of any choice of scale. FFGFT makes its fundamental
    statements on this layer.

  \item[Layer 2 --- the conversion to SI.] The moment one asks how many
    kilograms a certain mass has, how many seconds a certain duration
    lasts, or how many volts a certain potential carries, one enters a
    second layer: the translation between the dimensionless ratios of
    Layer 1 and the humanly chosen units kilogram, second, meter, volt.
    This translation requires a \emph{scale-bridge}, and in this bridge
    --- and only in this bridge --- the fractal corrections sit.
\end{description}

This is the central, often overlooked statement:

\begin{center}
\fbox{\parbox{0.85\textwidth}{\centering
  \textbf{Fractal geometry sits in the ratios themselves ---
  but there as a closed, frozen recursion.\\[4pt]
  It becomes visibly active only in the bridge to the SI language,
  where the embedding correction must be made explicit.}
}}
\end{center}

The ratios of Layer~1 are not fractal-free --- they \emph{are} fixed
points of fractal recursions that have come to rest. What is absent is
not the fractality but its visible motion. Only when one asks
\emph{how many kilograms} a mass has does the fixed point break open,
and the embedding correction must be carried along explicitly. This
transition --- not the physics itself, but the representation of
physics in human units --- pays the ``price'' that Doc.~185
(\emph{Time as the embedding cost of mathematics}) speaks of.

\section{Where the recursion is hidden}

At this point a misunderstanding can arise. If Layer 1 contains the
pure ratios, sharp and without correction --- where then is the
recursion of which FFGFT, as a \emph{fractal-geometric} theory, speaks?
Has one lost it among the ratios?

The answer is: no. It is preserved in the ratios. But there it looks
still.

A plain number like the ratio $m_{\mu}/m_{e} \approx 206{.}768$ looks
static. But it \emph{is} the endpoint of a recursion that has come to
rest in this number. Just as the golden section
$\varphi = (1+\sqrt{5})/2$ apparently is a simple number but in truth
the limit of the Fibonacci sequence $a_{n+1} = a_{n} + a_{n-1}$ ---
the recursion does not sit beside $\varphi$, it \emph{is} $\varphi$.
Just as $\pi$ sits in the $n$-gons and $e$ in the compound-interest
formula: a pure number is always a frozen process. A \emph{fixed
point} is a geometry that has absorbed its own motion.

The FFGFT ratios are no different. The lepton ratio $m_{\mu}/m_{e}$ is
the fixed point of a winding recursion on the 4D-torus, which converges
after three generations with step size $\xi^{1/3}$. The rational
prefactors $4/3, 16/5, 25/9$ (Doc.~205) are not arbitrary numbers but
standing waves of a recursive mode structure. The deeper mechanism
behind the choice of just these values is a Fibonacci-like convergence:
the recursion amplitude $\varphi^{-2n}/\sqrt{5}$ reaches the order of
$\xi_{0}$ at $n \approx 8$, and exactly there the mode sequence closes
into a finite table (Docs.~011, 159). Three generations of step size
$\xi^{1/3}$, terminated at Fibonacci convergence depth $n \approx 8$
--- this is not coincidence but geometric necessity that only eight
levels deep \emph{can} be reckoned before the convergence amplitude
drops below $\xi$.

The recursion is thus fully present in the ratios --- but it has been
put to rest because it has converged. What we see as a ``ratio'' is
the final form of a motion that has come to rest. Whoever sees a ratio
sees a finished story. Whoever asks \emph{why exactly this} ratio,
asks after the recursion that led to this fixed point. FFGFT answers
on both levels: the ratios are the results, the windings are the
underlying story.

\section{Where the recursion becomes visible}

If the ratios are closed in themselves --- where then does the
recursion become visible \emph{in motion}? Answer: in the bridge to
the SI language.

An SI unit (kilogram, second, meter) is not a fixed point; it is an
\emph{arbitrarily chosen intermediate station} on the scale ladder of
nature. The Planck scale and the electroweak scale lie about 17 orders
of magnitude apart; the kilogram sits somewhere in the middle, fixed
by a historically grown choice of scale, not by a geometric eigenvalue.
Whoever translates from a geometric eigenvalue in Layer 1 to a
kilogram value in Layer 2 does not fall from one finished recursion
into another finished one, but moves along a scale staircase that is
at this moment still \emph{in motion} --- not yet converged.

It is exactly there that $K_{\mathrm{frak}} = 1 - 100\xi$ appears. The
factor 100 is not ad hoc but the cumulative effect of the minimal
deviation $\xi$ over all rungs between Planck and electroweak scale.
The logarithmic distance $\ln(10^{17}) \approx 39$, geometrically
averaged over the relevant rungs, yields the amplification factor
$\sim 100$. It is a recursion that has \emph{not yet come to rest},
because our choice of scale cuts it off in the middle (Doc.~133).

The fractal geometry thus shows itself in two ways:

\begin{itemize}\setlength\itemsep{2pt}
  \item in Layer 1 as \emph{recursion at rest} ---
        in the ratios that are fixed points of a converged
        winding sequence;
  \item in the bridge as \emph{recursion in motion} ---
        in $K_{\mathrm{frak}}$, the cumulative trace of a
        scale cascade that we interrupt mid-way by our choice
        of units.
\end{itemize}

\noindent
Both are the same geometry. Once as fixed point, once as scale flow.
What looks like ``two different phenomena'' is the same fractal
self-similarity in two states: at the goal and on the way.

\section{Consequences}

From this architecture four consequences follow that noticeably
simplify the understanding of many individual FFGFT results.

\subsection{Ratio predictions are sharp}

When FFGFT makes a statement about a \emph{ratio} --- for instance
about $m_{\mu}/m_{e}$ or about the muon-to-tauon lifetime ratio
(Doc.~160) ---, this prediction is sharp: \emph{without} fractal
correction, because the correction cancels in a ratio of quantities of
the same kind.

\begin{center}
\fbox{\parbox{0.85\textwidth}{
  When both quantities in numerator and denominator pass to the SI
  scale through the same conversion factor, this factor cancels out.
  What remains is pure geometry.
}}
\end{center}

For this reason FFGFT ratio predictions are typically accurate to
four or five decimal places. They test \emph{the geometry}, not the
conversion.

\subsection{SI values carry the fractal signature}

The moment FFGFT makes a statement about an \emph{absolute SI value}
--- for instance about the electron mass in kilograms, or about the
gravitational constant in m\textsuperscript{3}/(kg$\cdot$s\textsuperscript{2})
---, the fractal correction appears. This is not a deficit of the
theory but a necessary consequence of the architecture:

\begin{itemize}\setlength\itemsep{0pt}
  \item In Layer 1 the electron mass is the geometric number
        $m_{e}^{\,\mathrm{T0}} = 0{.}511$ (Doc.~014).
  \item The conversion to kilograms requires the scale factor
        $S_{T0} \approx 1{.}782662\times 10^{-30}$, which is not chosen
        arbitrarily but derived from the fractal geometry.
  \item Only the combination of both --- geometric number times scale
        factor --- yields the experimental value
        $m_{e}^{\,\mathrm{SI}} \approx 9{.}1094 \times 10^{-31}$~kg.
\end{itemize}

The fractal correction $K_{\mathrm{frak}} = 1 - 100\xi \approx 0{.}9867$
(Doc.~133) belongs in this bridge, not in the electron mass itself.

\subsection{$D_f$ is a property of the bridge}

A clarification on the fractal dimension $D_{f} \approx 2{.}999867$
follows: it does not sit ``in the mass'', nor ``in space'' in the
naive sense, but in the bridge between geometrically natural
language and SI language.

\begin{itemize}\setlength\itemsep{0pt}
  \item The \emph{local} fractal dimension of spacetime
        $D_{f}^{\,\text{space}} = 3 - \xi \approx 2{.}9999$ is a
        property of the vacuum and supplies only corrections of order
        $10^{-5}$, that is, practically one.
  \item The \emph{effective} fractal dimension
        $D_{f}^{\,\text{eff}} \approx 2{.}973$ that appears in the
        bridge is a \emph{cumulative} quantity: it arises from the
        effect of the local deviation $\xi$ over the full
        renormalization-group flow from the Planck scale to the
        electroweak scale (Doc.~133).
\end{itemize}

Both have the same geometric origin; but they sit at different
positions of the architecture, and they must not be confused.

\subsection{$K_{\mathrm{frak}}$ is not a fitting constant}

A common question to FFGFT runs: ``Is $K_{\mathrm{frak}} = 1 - 100\xi$
not simply chosen so that the experimental values come out?'' The
answer becomes clear through the two-layer architecture:

\begin{itemize}\setlength\itemsep{0pt}
  \item $K_{\mathrm{frak}}$ is \emph{not} a free constant. It is the
        only geometrically possible bridge function between Layer 1
        and Layer 2, given by the cumulative effect of $\xi$ over
        the RG flow.
  \item The factor $100$ is not ad hoc but follows from the
        logarithmic distance between Planck and electroweak scale
        ($\ln(10^{17}) \approx 39$) and the geometric averaging over
        the relevant rungs (Doc.~133).
  \item Once $\xi = 4/30000$ is fixed, $K_{\mathrm{frak}}$ is
        \emph{compulsorily} determined.
\end{itemize}

There is therefore still only one free parameter: $\xi_{0}$.
Everything that appears in Layer 2 follows from this one parameter
and from the geometry of the scale bridge.

\section{The picture}

Put into words, the two-layer architecture says:

\begin{center}
\fbox{\parbox{0.85\textwidth}{\centering
  \textbf{Nature reckons in ratios.}\\[2pt]
  \textbf{We measure in SI.}\\[2pt]
  \textbf{The fractal geometry sits in the translator.}
}}
\end{center}

This is not only a pedagogical rule of thumb but a structural statement
about FFGFT. It explains why some predictions are exact and others
receive a $1{.}3\%$ correction; it explains why $K_{\mathrm{frak}}$ is
not a fitting parameter; it explains why the theory operates with one
free parameter; and it makes precise the relation between geometric
idealization and experimental measurement.

The epistemic point runs: \emph{what we call ``absolute values in
SI'' are human artifacts}. The universe ``knows'' nothing of meter,
second, or kilogram. It knows only ratios, winding numbers, phases.
The fractal corrections are the price we pay for measuring the
universe \emph{through our arbitrary units}. They are not a defect of
the world but a property of the language in which we describe the
world.

\section{The deeper point: the universe does not compute}

Even the formulation ``Layer 1 \emph{computes} ratios, Layer 2
\emph{computes} the SI values'' is still too mechanistic. In neither
layer does any computation take place. The universe does not compute;
it \emph{is} the configuration.

Layer 1 is a \emph{geometry of admissible mode configurations}. A
mode either exists or does not exist. It is not ``computed''; it is
the geometric configuration itself. The ratios stand where they stand
because the winding recursion has converged at that point --- not
because anyone has computed it.

Layer 2, the bridge, is likewise not a computation but a
\emph{translation}. $K_{\mathrm{frak}}$ is not computed but is the
only geometrically possible bridge function. The scale staircase
itself does not ``compute''; it is there.

Computation takes place only where an observer wants to derive Layer 2
from Layer 1 --- the computation happens in the \emph{observer}, not
in the world.

The same logic makes photonic computing media (TFLN crystals, Ising
machines, analog resonators, quantum dots; see Docs.~161, 167, 173)
so efficient: the medium does not ``compute'', it lies in the
solution. The geometry of the admissible phase configurations
\emph{is} the result. What in a digital machine runs as a sequential
operation in time is in the photonic medium a \emph{simultaneous}
configuration in space: not algorithm but geometry. FFGFT makes the
same claim about nature as a whole --- the mode structure of the
4D-torus is not the program that nature executes, but nature itself.

\section{Where we really are: $c$ as the experience of transience}

But the architecture is still not complete. Layer 1 is the geometric
pure form; Layer 2 is the human scale description. Both are
\emph{languages} about something in which we ourselves actually are.
But we live neither in the pure geometry nor in the SI values. We
live in a third position that does not coincide with either layer.

What we actually experience is \emph{transience}. For us $c$ is not
one. $c$ has a finite value because we experience the light path as
duration, as a distance that takes time. Not because $c$ ``really''
were $299\,792\,458$~m/s --- in pure geometry $c = 1$ ---, but
because we sit inside the configuration and experience it as
\emph{passing}. Were we Layer 1 itself, we would no longer be
observers with a lifespan, but the geometry. The values set to one
belong to the eye of the geometry. We see from inside --- and
exactly there they take finite values, and exactly there they are
inserted, as soon as we measure.

This is the sense of the statement that the values set to one must be
reinserted when projected measurement values are generated: not
because the geometry has them, but because \emph{we who measure} are
in a position where they take finite values. $c$ is the
phenomenological trace of our inside position. Without an observer
who ages, there would be no finite $c$.

\section{Closing point: the fractal trace of our experience}

And with this the architecture, which at the beginning looked merely
static, closes into a dynamic form. The recursion was never only in
the ratios (Layer 1) or only in the bridge (Layer 2). It is also in us.

Every moment in which we age contains all previous moments in damped
form. Our experience is not a linear stream but a \emph{concert hall}
(Doc.~159): every new sound carries the echo of all old ones. Human
memory is not a storage but a recursive resonance in which every
recurrence is a faint copy of the previous one.

If the universe \emph{is} the configuration and we \emph{experience}
the configuration, then our experience is the \emph{fractal trace of
the geometry in which we lie}. The recursion is not an abstract
feature of the theory --- it is what we recognize ourselves in as
observers. The same pattern that has come to rest in the ratios and
that is still in motion in the scale bridge is also the pattern of
our perception.

\medskip

This is the arc this document draws: from the apparent multiplicity
of natural constants over the two layers and their recursive
relation to the position of the observer, who does not \emph{know}
the geometry but \emph{is} it --- and yet runs through it as
transience. FFGFT does not describe a machinery that someone looks at
from outside. It describes a geometry that carries itself, and an
observer who lies inside this geometry and must therefore express it
in two languages.

Whoever has seen this no longer sees $c$, $\hbar$, $\alpha$, and
$K_{\mathrm{frak}}$ as separate natural constants but as different
traces of the same one geometry --- seen from different positions
within the one configuration that we do not observe but in which we are.

\section{Reading order: where each layer is worked out}

For systematic study, the two layers are treated in the following
documents:

\begin{itemize}\setlength\itemsep{0pt}
  \item \textbf{Layer 1 (ratios, geometric pure form):}
    Doc.~002 (Journey), Doc.~015 (systematics of natural units),
    Doc.~009 (origin of $\xi$), Doc.~205 (narrative),
    Doc.~006 (particle masses as ratios),
    Doc.~060/116/159 (mode and winding ratios),
    Doc.~230 (Hilbert-space bridge, likewise ratio-based).
  \item \textbf{Layer 2 (scale bridge, fractal corrections):}
    Doc.~014 (transition nat./SI with scale factor $S_{T0}$),
    Doc.~013 (SI conventions), Doc.~133 (complete derivation of
    $K_{\mathrm{frak}}$), Doc.~134 (c as convention),
    Doc.~044 (conventions for $\alpha$), Doc.~190 (corrections
    register), Doc.~185 (embedding cost of mathematics),
    Doc.~210 (winding corrections).
\end{itemize}

Whoever reads the documents in this order sees the architecture this
document speaks of at concrete points: ratios above, bridge below, no
free parameter except $\xi_{0}$.

\section*{Summary}

\begin{itemize}\setlength\itemsep{2pt}
  \item FFGFT operates on two cleanly separated layers:
        ratios (Layer 1, fractal-free) and SI projection (Layer 2,
        with fractal bridge).
  \item The values $c, \hbar, G, k_B, \alpha$ that have been ``set to
        one'' do not disappear but are factored out; they return when
        projected measurement values are generated.
  \item The recursion is present in the ratios as recursion at rest
        (fixed points of a winding sequence with Fibonacci convergence
        at $n \approx 8$) and in the bridge as recursion in motion
        ($K_{\mathrm{frak}} = 1 - 100\xi$).
  \item The universe does not compute; it is the configuration.
        Computation happens in the observer --- analogous to the logic
        of photonic computing media (TFLN, Ising machines).
  \item We live neither in Layer 1 nor in Layer 2 but in a third
        position: the experience of transience. $c$ has for us a
        finite value because we \emph{traverse} the configuration
        instead of \emph{being} it.
  \item The recursion also shows itself in our experience: every
        moment carries the echo of all previous ones. Our perception
        is the fractal trace of the geometry in which we lie.
\end{itemize}

\medskip
\noindent\textit{Status: May 2026. This document is methodologically
pedagogical and contains no independent physical claims; all factual
statements refer to the source documents cited.}
